\section{Implementation Details}
\label{app:implementation}
Here we introduce the implementation of each unlearning method used in~\cref{sec:experiment}, including hyperparameters, training protocols, and compute budgets. All methods are applied to CompVis Stable Diffusion v1.4~\citep{rombach2022high}. The full implementation, including code, prompt CSVs, and evaluation scripts, is available at our public repository:
\url{https://github.com/Jazzssmine/concept-entangle}

\subsection{Method-Specific Hyperparameters}
We group the 13 benchmarked methods into three mechanism families.
\emph{Closed-form editing methods} (UCE, RECE) modify model weights through analytic updates. 
\emph{Fine-tuning methods} (ESD, SalUn, EDiff, CPE, TRUST, AdvUnlearn, STEREO) modify model weights through gradient-based optimization.
\emph{Inference-time methods} (SEOT, SAeUron, SLD, SAFREE) intervene during generation without permanently modifying the base model. AdvUnlearn and STEREO additionally incorporate adversarial training, explicitly targeting the aggressive-erasure end of the trade-off; their inclusion tests whether robustness gains incur measurable cost to neighbor preservation.

Unless otherwise noted, we use the hyperparameters specified in each method's original publication. We document deviations and choices not fully specified by the original papers in the per-method descriptions below.

\paragraph{UCE~\citep{gandikota2024unified}.}
Unified Concept Editing uses closed-form updates to text-to-image attention projections to erase target concepts without gradient training. Configuration: guide concept \textbf{unconditional}, erase scale \textbf{1.0}, preserve scale \textbf{1.0}, regularization $\lambda$ \textbf{0.5}, no explicit preserve concepts; each concept is erased independently.

\paragraph{RECE~\citep{gong2024reliable}.}
Reliable and Efficient Concept Erasure iteratively discovers embeddings that recover the target concept and removes their effect through closed-form cross-attention edits.
Configuration: initialized from base SD v1.4, edit epochs \textbf{1}, erase scale \textbf{1.0}, preserve scale \textbf{0.1}, $\lambda$ \textbf{0.1}, adversarial-embedding regularization \textbf{0.1}, guide concept \textbf{unconditional}.


\paragraph{ESD~\citep{gandikota2023erasing}.} 
Erased Stable Diffusion fine-tunes the cross-attention layers to suppress the target concept via a negative-guidance objective. We use the ESD-sd (stable-diffusion) variant. 
Configuration: learning rate \textbf{5e-5}, training steps \textbf{200}, negative guidance scale $\eta = \textbf{5}$, 
batch size \textbf{1}. Training uses the target concept name as the single forget prompt.
For the cross-architecture experiments (\cref{app:arch}), we adapt ESD to each backbone.
On SDXL (base-1.0) we use the ESD-x variant, which fine-tunes only the U-Net cross-attention layers (\texttt{attn2} query, key, value, and output projections; about 546M parameters). Configuration: Adam, learning rate \textbf{2e-4}, training steps \textbf{200}, batch size \textbf{1}, negative guidance scale $\eta = \textbf{1}$.
On FLUX.1-schnell we use a stricter ESD-x variant that fine-tunes only the attention key and value projections (\texttt{to\_k}, \texttt{to\_v}) in all 19 joint and 38 single transformer blocks (about 1.08B parameters), keeping the text-side projections (\texttt{add\_k\_proj}, \texttt{add\_v\_proj}) frozen. Configuration: learning rate \textbf{1e-5}, training steps \textbf{500}, negative guidance scale $\eta = \textbf{2}$.
% Provenance (not for the paper): SDXL settings come from the saved run configs for dog and castle; the SDXL cat run has no saved config and is assumed to match. FLUX lr and guidance are taken from the run folder name *_strict_lr1e5_ng2.

\paragraph{SalUn~\citep{fan2023salun}.} 
Saliency-guided Unlearning identifies salient weights via gradient magnitudes and applies random labeling on the forget set to those weights only. We use saliency threshold sparsity \textbf{0.5}, forget-set size \textbf{800}, learning rate \textbf{1e-5}, and training epochs \textbf{5}. The forget and retain sets consist of 800 target and 800 non-target images generated by Stable Diffusion v1.4.
\paragraph{EDiff~\citep{wu2024erasediff}.} 
EraseDiff fine-tunes the UNet using a bi-level optimization that explicitly balances erasure and retention objectives. Configuration: learning rate \textbf{1e-5}, inner steps per outer update \textbf{2}, total outer iterations \textbf{300}.
\paragraph{CPE~\citep{lee2025concept}}
Concept Pinpoint Eraser trains nonlinear residual attention gates with an anchoring loss to suppress target concepts while preserving other concepts.
Configuration: \textbf{5} stages of \textbf{450} iterations each (stage 1 uses $4\times$ the iterations and $10\times$ the learning rate), learning rate \textbf{3e-5}, gate rank \textbf{16}, erasure guidance \textbf{0.3}, PAL weight \textbf{1e5}, noise \textbf{1e-3}; adversarial learning with learning rate \textbf{1e-2}, \textbf{450} iterations, and \textbf{16} added prompts per stage.

\paragraph{TRUST~\citep{kori2026selective}}
TRUST dynamically identifies target-associated neurons and selectively fine-tunes them with Hessian-based regularization.
Configuration: concept neurons are heads whose $|z\text{-score}| > $ \textbf{2.0} on the CLIP-loss attribution; learning rate \textbf{1e-4}, batch size \textbf{2}, \textbf{2} epochs of \textbf{60} steps (final checkpoint), \textbf{25} sampling steps for neuron localization.

\paragraph{AdvUnlearn~\citep{zhang2024defensive}.} 
Adversarial Unlearning augments the erasure objective with adversarial prompt generation to improve robustness against indirect concept recovery. We use adversarial prompt budget 30 per training step, adversarial learning rate \textbf{1e-3}, main-model learning rate \textbf{1e-5}, and training steps \textbf{1000}. The attack prompts are updated per training steps.
\paragraph{STEREO~\citep{srivatsan2025stereo}.} 
STEREO proceeds in two phases: a Search for Trainable Embeddings via RObust optimization (STE) phase that learns adversarial token embeddings $v_1^*, v_2^*$, followed by a Robust Erasure via Orthogonalization (REO) phase that suppresses the subspace spanned by these embeddings. STE phase: \textbf{200} steps, learning rate 
\textbf{5e-6}. REO phase: \textbf{200} steps, learning rate \textbf{2e-5}, $K = \textbf{2}$ adversarial embeddings per target.
\paragraph{SEOT~\citep{li2024get}.} 
Soft Embedding Optimization for Target-concept erasure modifies the text embedding at inference time to suppress the target concept without altering model weights. We use embedding regularization strength \textbf{0.01} and optimization steps per generation \textbf{10}.
\paragraph{SAeUron~\citep{cywinski2025saeuron}.} 
SAeUron uses sparse autoencoder feature ablation applied at inference: for each target, $\tau_c$ concept-specific SAE features are identified from activation statistics, and these features are zeroed during generation. We use the pre-trained SAE released by the authors, with 
$\tau_c = \textbf{1}$ features ablated per target. 


\paragraph{SLD~\citep{schramowski2023safe}.}
Safe Latent Diffusion adds safety guidance during denoising to suppress target content without retraining the model.
Configuration: the \textbf{SLD-Max} preset, i.e.\ safety guidance scale \textbf{5000}, warm-up steps \textbf{0}, threshold \textbf{1.0}, momentum scale \textbf{0.5}, momentum $\beta$ \textbf{0.7}; \textbf{50} sampling steps, CFG scale \textbf{7.5}.

\paragraph{SAFREE~\citep{yoon2025safree}.}
SAFREE filters prompt embeddings away from a target-concept subspace and adapts its intervention during denoising without updating model weights.
Configuration: concept subspace from \textbf{12} per-concept negative terms (synonyms, sub-breeds, and phrasings such as ``a photo of a dog''), $\alpha$ \textbf{-0.5}, re-attention up to step \textbf{20}, FreeU hyperparameters \textbf{(1.0, 1.0, 0.9, 0.2)}; \textbf{50} sampling steps, CFG scale \textbf{7.5}.


\subsection{Compute and Infrastructure}
\label{app:compute}
All experiments were run on GPU A100 40GB. Unlearning per method-target pair requires approximately 1 GPU-hours for fine-tuning methods and 2 GPU-hours for inference-time methods. Evaluation (image generation + classification) requires approximately 0.5 GPU-hours per method-target pair.

\subsection{Reproducibility}
\label{app:reproducibility}
For each (method, target) pair, we generate 100 images from direct prompts, 100 to 235 from indirect prompts (depending on how many valid indirect prompts the target yields; 235 for \texttt{castle}, 200 for \texttt{dog}, 195 for \texttt{horse}, and 100 for every other target), 50 per neighbor candidate (500 total), and 50 per control concept (150 total); IRA is computed over the 650 neighbor and control images. All images use fixed random seeds, with seed equal to the row index in the target's prompt CSV. This ensures seed-level variation is controlled when comparing methods on the same prompt.
The base model $\mathcal{D}_\theta$ is loaded from the HuggingFace checkpoint \texttt{CompVis/stable-diffusion-v1-4}. Code, prompt CSVs, and the classification vocabulary (\cref{tab:vocab}) will be released to enable full reproduction of~\cref{tab:main_results_avg}.

\paragraph{Use of existing assets.}
Our experiments use publicly released assets from prior work: the CompVis Stable Diffusion v1.4 checkpoint, released under the CreativeML OpenRAIL-M license; LAION metadata/captions, released under CC-BY 4.0 while the underlying images remain subject to their original copyrights; and the released SAeUron sparse-autoencoder features, released under Apache-2.0. We cite the original sources for these assets and use them only for research evaluation under their stated release terms and licenses.